% The viewers on their own: what each of the three viewers is for, installing and deploying them
% without Medorama, their security model (the root, the settings that narrow it, symbolic links,
% hidden entries, the shared password), the ports, the themes, the website and its examples.
% Sources: viewers/README.md, docs/building.md ("The build options of a viewer alone"), checked
% against viewers/vitalstream/rust/src/{main,settings,server,auth,auth_none}.rs and
% viewers/ecg/python/muse/muse_reader_py/page_server.py (the same file as dicom-py's).
\chapter{The viewers on their own}
\label{chap:viewers-alone}

Medorama shows a patient's recordings, ECGs and images in three viewers: VitalStream, the ECG
viewer and the PACS viewer. Each is a program of its own, in its own folder of the repository
(\file{viewers/vitalstream}, \file{viewers/ecg}, \file{viewers/pacs}), with its own design, and
each is useful without Medorama, because it opens the files given to it. Medorama serves them and
opens them as tabs of a patient's page, at an address that names what to show
(\cref{chap:viewers-in-medorama}).

This part of the manual describes the viewers alone, for researchers who use them on their own
computers and administrators who deploy them. This chapter describes what they have in common:
what each is for, how to install and deploy it on one computer or on a server, how it decides
which files it may open, and what Medorama adds when it runs them. The three chapters that follow
describe each viewer in full: VitalStream (\cref{chap:vitalstream}), the ECG viewer
(\cref{chap:ecg-viewer}) and the PACS viewer (\cref{chap:pacs-viewer}).

% ==========================================================================================
\section{The three viewers}
\label{sec:viewers-alone-three}

\begin{xltabular}{\linewidth}{@{}P{2.1cm}LLP{3.3cm}@{}}
\caption{The three viewers: what they show, what they open, and how they run alone.}
\label{tab:viewers-alone-three}\\
\toprule
{\tabhead Viewer} & {\tabhead Shows} & {\tabhead Opens} & {\tabhead Alone} \\
\midrule
\endfirsthead
\toprule
{\tabhead Viewer} & {\tabhead Shows} & {\tabhead Opens} & {\tabhead Alone} \\
\midrule
\endhead
\bottomrule
\endfoot
VitalStream & Waveforms: ICU monitoring (every channel, the vital signs, the alarms), Holter
  ECGs, EEGs and sleep studies of hours, days or weeks, on one time axis, from weeks down to
  single samples & WFDB records (single- and multi-segment), folders of records as one session,
  WFDB collections; EDF, EDF+, BDF and BDF+ files; other formats (such as \file{.Stp} files)
  through an installed adapter & A server: \cmd{vitalstream-cpp}, \cmd{vitalstream-rs} or
  \cmd{vitalstream-py} \\
The ECG viewer & 12-lead ECGs on standard paper (25\,mm/s, 10\,mm/mV, a 1\,mm and 5\,mm grid):
  3~$\times$~4 with a rhythm strip, 6~$\times$~2 with a rhythm strip, 12~$\times$~1; calipers &
  WFDB records (single- and multi-segment), GE MUSE XML, DICOM ECGs, EDF and BDF files, the
  viewer's JSON form, PNG, JPEG, PDF; a folder of ECGs, one at a time & A page:
  \file{viewers/ecg/index.html} from disk, on the website, or served by \cmd{muse-py serve} \\
The PACS viewer & Cardiology and other DICOM studies: echo and angiography loops at their frame
  rate, stacks of slices (CT, MR, NM, X-ray) with window and level, structured reports, PDF
  documents; calipers in cm or mm & DICOM (native, RLE, JPEG baseline, extended and lossless),
  MP4, WebM, MOV, PNG, JPEG, PDF; a folder of studies & A page: \file{viewers/pacs/index.html}
  from disk, on the website, or served by \cmd{dicom-py serve} \\
\end{xltabular}

\paragraph{VitalStream} is a web server and its page. It reads the recordings where they lie, on
the computer it runs on, indexes each recording once, and sends the page only the samples that
the view needs. Long, high-frequency recordings are its strength: a day of ICU data, a week of
Holter ECG or a night's sleep study opens at once and stays fluid while zooming. It adds nothing
that is not in the recording. Samples are never resampled, interpolated or stitched, and where
the recording is ambiguous (the length of a break, a record without a start time) the page says
so. It comes in
three implementations, in Python, Rust and C++, that serve the same page, answer the same
requests and write the same index (\cref{chap:vitalstream}).

\paragraph{The ECG and PACS viewers} are static pages: HTML, CSS and plain scripts, with no
dependencies and no build step. They work from disk (\code{file://}) and from any web server.
Everything is read and drawn in the browser, so the files a user picks are read on the user's own
computer and sent nowhere. Their Python packages, \cmd{muse-py} and \cmd{dicom-py}, add a
small server, \cmd{serve}, that serves the page together with the folders of the computer it
runs on, so that the ECGs or studies of a server are read without copying them
(\cref{chap:ecg-viewer,chap:pacs-viewer}).

% ==========================================================================================
\section{Trying them}
\label{sec:viewers-alone-try}

The ECG and PACS viewers run on Medorama's website, in the browser. Each has a page of synthetic
examples (every value and every image is invented, and no file describes a real person), where
each file has a link that opens it in the live viewer.

\begin{table}[htbp]
\centering
\caption{The live viewers and their examples.}
\label{tab:viewers-alone-try}
\begin{tabularx}{\linewidth}{@{}P{2.1cm}P{5.6cm}L@{}}
\toprule
\tabhead Viewer & \tabhead Live & \tabhead Examples \\
\midrule
ECG & \url{https://medorama.org/viewers/ecg/} & \file{viewers/ecg/examples/}:
  a WFDB record, a GE MUSE XML file, a DICOM ECG, a PDF and a PNG printout
  (\cref{chap:ecg-viewer}) \\
PACS & \url{https://medorama.org/viewers/pacs/} &
  \file{viewers/pacs/examples/}: an echo loop, a coronary angiography run, a chest X-ray
  (\cref{chap:pacs-viewer}) \\
VitalStream & None: it is a server, to install (\cref{sec:viewers-alone-install}) &
  \file{viewers/vitalstream/testdata/}: a WFDB collection, plain records, a folder of records
  to open as one session, and EDF files (\cref{chap:vitalstream}) \\
\bottomrule
\end{tabularx}
\end{table}

The website's tabs \ui{Viewers}, \ui{VitalStream}, \ui{ECG} and \ui{PACS} lead to the viewers'
documents; the examples pages are at \url{https://medorama.org/viewers/ecg/examples/}
and \url{https://medorama.org/viewers/pacs/examples/}. The ECG and PACS
examples are written by \file{tools/make-viewer-examples.py}, VitalStream's by
\file{viewers/vitalstream/testdata/generate.py}. From a clone of the repository, open
\file{viewers/ecg/index.html} or \file{viewers/pacs/index.html} in a browser, then
\ui{Open multiple records\dots} and choose the folder \file{examples/}.

% ==========================================================================================
\section{Installing a viewer alone}
\label{sec:viewers-alone-install}

\Cref{tab:viewers-alone-install} sums up how each viewer is installed. The website gives the
commands for Linux, macOS and Windows, one after the other, at \url{https://medorama.org/docs/install.html}.

\begin{table}[htbp]
\centering
\caption{Installing each viewer without Medorama.}
\label{tab:viewers-alone-install}
\begin{tabularx}{\linewidth}{@{}P{2.4cm}L@{}}
\toprule
\tabhead Viewer & \tabhead Install \\
\midrule
VitalStream & \cmd{viewers/vitalstream/install.sh} (Windows: \file{install.ps1}) builds one
  viewer, with a root it never leaves (by default the home folder) and optional password
  support, into \file{\textasciitilde/.local/bin} (Windows:
  \code{\%LOCALAPPDATA\%\textbackslash vitalstream\textbackslash bin}) (\cref{chap:vitalstream}) \\
The ECG viewer & Nothing: open the page. To serve it from a computer:
  \cmd{pip install ./viewers/ecg/python/muse}, then \cmd{muse-py serve} \\
The PACS viewer & Nothing: open the page. To serve it from a computer:
  \cmd{pip install ./viewers/pacs/python/dicom}, then \cmd{dicom-py serve} \\
\bottomrule
\end{tabularx}
\end{table}

From a clone of the repository:

\begin{lstlisting}[style=shell]
git clone https://github.com/alphanumericslab/medorama.git && cd medorama
viewers/vitalstream/install.sh                    # VitalStream (C++), into ~/.local/bin
viewers/vitalstream/install.sh --root /data/icu   # one that opens nothing outside /data/icu
python3 -m venv .venv
.venv/bin/pip install ./viewers/ecg/python/muse ./viewers/pacs/python/dicom   # muse-py, dicom-py
\end{lstlisting}

The prerequisites of each build (a C++20 compiler, Rust, CMake, Python) and the builds on
machines without administrator rights or without internet access are described in
\cref{chap:installation}.

\subsection{The build options of a viewer alone}
\label{sec:viewers-alone-build-options}

A viewer installed on its own has two options fixed when it is built: its \emph{root}, the
folder it never opens anything outside of, and whether it supports a password. The installers
pass them; a build by hand sets them as in \cref{tab:viewers-alone-build-options}. For the root,
\code{\textasciitilde} stands for the home folder of whoever runs the viewer, and \code{/} for the
whole file system; without the option, the root is the whole file system.

\begin{xltabular}{\linewidth}{@{}P{2.9cm}LP{3.6cm}@{}}
\caption{The build options of a viewer alone.}
\label{tab:viewers-alone-build-options}\\
\toprule
{\tabhead Program} & {\tabhead Root} & {\tabhead Without password support} \\
\midrule
\endfirsthead
\toprule
{\tabhead Program} & {\tabhead Root} & {\tabhead Without password support} \\
\midrule
\endhead
\bottomrule
\endfoot
\cmd{vitalstream-rs} & \code{VITALSTREAM\_ROOT=DIR cargo build -{}-release -p vitalstream-rs
  -{}-target-dir target/standalone} & \opt{-{}-no-default-features -{}-features wfdb} (the Cargo
  feature \code{password} left out) \\
\cmd{vitalstream-cpp} & \code{cmake -{}-preset viewer -B build/standalone -DVS\_ROOT=DIR} &
  \code{-DVS\_PASSWORD=OFF} \\
\cmd{vitalstream-py} & \code{VITALSTREAM\_ROOT=DIR pip install
  './viewers/vitalstream/python/vitalstream[wfdb]'} & \code{VITALSTREAM\_PASSWORD=0} \\
\cmd{muse-py serve}, \cmd{dicom-py serve} & \code{VIEWER\_ROOT=DIR pip install
  ./viewers/ecg/python/muse} (or \file{./viewers/pacs/python/dicom}) & \code{VIEWER\_PASSWORD=0} \\
\end{xltabular}

The standalone installer builds in folders of its own (\file{target/standalone},
\file{viewers/vitalstream/cpp/build/standalone}), so that the viewer Medorama runs, built in
\file{target/release} or \file{build/viewer}, keeps its own options. Medorama's installer,
\cmd{tools/install-medorama.sh}, builds its viewer with the whole file system as its root and
with password support, since Medorama applies its own rules to what each user opens
(\cref{chap:viewers-in-medorama}). WFDB support is a third build option, for every program
(\cref{sec:installation-no-wfdb}).

% ==========================================================================================
\section{Deploying}
\label{sec:viewers-alone-deploy}

\subsection{On one computer}

Each viewer serves its page to the same computer, and opens the browser at its address:

\begin{lstlisting}[style=shell]
vitalstream-cpp            # http://127.0.0.1:8070 (vitalstream-rs: 8060, vitalstream-py: 8050)
muse-py serve              # http://127.0.0.1:8080/browse: the ECG viewer and the folders under ~
dicom-py serve             # http://127.0.0.1:8081/browse: the PACS viewer and the folders under ~
\end{lstlisting}

The ECG and PACS viewers need no server at all on one computer: their pages, opened from disk or
from the website, read the files picked in the browser. \cmd{serve} is for files that lie on
another computer.

\subsection{On a server, through an SSH tunnel}
\label{sec:viewers-alone-server}

On a server where the data are, nothing needs to be copied anywhere. Every viewer listens on
127.0.0.1 by default (that machine only); from another computer, an SSH tunnel carries its port:
\begin{enumerate}
\item Install the viewer on the server, with a root that holds the data
  (\cref{sec:viewers-alone-build-options}), and set a password once
  (\cref{sec:viewers-alone-password}).
\item Start it without a browser, in \cmd{tmux} or as a service (\cref{chap:vitalstream}):
  \cmd{vitalstream-cpp -{}-no-browser}, \cmd{muse-py serve -{}-root /data/ecg -{}-no-browser},
  \cmd{dicom-py serve -{}-root /data/echo -{}-no-browser}.
\item On your own computer, open a tunnel to its port, and open the address:
\end{enumerate}

\begin{lstlisting}[style=shell]
ssh -N -L 8070:127.0.0.1:8070 user@server     # VitalStream (C++): http://127.0.0.1:8070
ssh -N -L 8080:127.0.0.1:8080 user@server     # muse-py serve: http://127.0.0.1:8080/
ssh -N -L 8081:127.0.0.1:8081 user@server     # dicom-py serve: http://127.0.0.1:8081/
\end{lstlisting}

Without a password, anyone who can connect to the port can read the files under the viewer's
root: on a machine that others can log in to, set one. The pages use relative addresses, so
VitalStream also works behind a portal or a reverse proxy that serves it under a path prefix.
Each viewer answers only requests that name it by a loopback name, an IP address or its own
host, which defeats DNS rebinding: a VitalStream viewer reached by another name (\code{host =
"0.0.0.0"}, or a proxy that passes its public name on) lists that name in \code{allowed\_hosts},
and behind an HTTPS proxy sets \code{secure\_cookie = true} (\cref{sec:viewer-api-checks}).

\subsection{The ports}
\label{sec:viewers-alone-ports}

\begin{table}[htbp]
\centering
\caption{The default ports of the viewers and of Medorama.}
\label{tab:viewers-alone-ports}
\begin{tabularx}{\linewidth}{@{}LP{6.4cm}@{}}
\toprule
\tabhead Program & \tabhead Port (default) \\
\midrule
\cmd{vitalstream-py}, \cmd{vitalstream-rs}, \cmd{vitalstream-cpp} & 8050, 8060, 8070 \\
\cmd{muse-py serve} (ECG), \cmd{dicom-py serve} (PACS) & 8080, 8081 \\
Medorama (\cmd{medorama serve}) & 8090; it starts its own VitalStream viewer, on a free port \\
\bottomrule
\end{tabularx}
\end{table}

\opt{-{}-host} and \opt{-{}-port} change them (VitalStream: also the settings \code{host} and
\code{port}). Several viewers, and the three VitalStream viewers, can run side by side.

% ==========================================================================================
\section{The security model}
\label{sec:viewers-alone-security}

A viewer alone has no users, no database, no sign-in beyond an optional password, and no
de-identification: those belong to Medorama (\cref{sec:viewers-alone-medorama}). What protects
the files is the viewer's root, the folders its settings allow within it, and the password.

\subsection{The root}
\label{sec:viewers-alone-root}

\paragraph{Fixed when the viewer is built or installed.} The \emph{root} of a viewer is the folder
it never lists or opens anything outside of. It is set when the viewer is built (VitalStream:
\opt{-{}-root DIR} of \cmd{install.sh}, \code{VITALSTREAM\_ROOT}, \code{VS\_ROOT}) or when the
package is installed (\cmd{muse-py}, \cmd{dicom-py}: \code{VIEWER\_ROOT}), and cannot be widened
afterwards. \code{\textasciitilde} is the home folder of whoever runs the viewer. The standalone
VitalStream installer's default is \code{\textasciitilde}; a build without the option, such as a
developer's \cmd{cargo build}, \cmd{cmake -{}-preset release} or \cmd{pip install}, has the whole
file system as its root, and so has a \cmd{muse-py} or \cmd{dicom-py} installed without
\code{VIEWER\_ROOT}.

\paragraph{Narrowed by the settings, never widened.} At run time, the folders the viewer may
open are narrowed further:
\begin{itemize}
\item VitalStream: the setting \code{data\_roots}, or its option \opt{-{}-root DIR} (repeatable),
  names the folders it may open; by default, its root. A folder that contains the root is
  narrowed to the root; one outside it is left out, with a warning when the viewer starts
  (\code{DIR is outside this viewer's root (ROOT): left out}). When none is left, the viewer does
  not start (\code{no data root inside this viewer's root, ROOT (set when it was built)}).
\item \cmd{muse-py serve} and \cmd{dicom-py serve}: \opt{-{}-root DIR} names the one folder
  served; by default, the home folder. A folder that contains the package's root is narrowed to
  it; a folder outside it, or that is not a folder, is refused, and \cmd{serve} does not start
  (\code{DIR is not a folder inside this viewer's root, ROOT (set when it was installed)}). With
  a \code{VIEWER\_ROOT} that does not contain the home folder, \opt{-{}-root} must therefore be
  given.
\end{itemize}
The install's root and the viewer's \opt{-{}-root} are two different things: the first is fixed
in the program, the second narrows it each time the viewer starts.

\paragraph{Paths are checked after their symbolic links are resolved.} A path is allowed when the
file or folder it leads to, every symbolic link resolved, lies inside an allowed folder. So a
symbolic link inside the root that leads outside it is refused (VitalStream answers
\code{403}, \code{path outside the configured data roots}; \cmd{serve} answers \code{403}), and
it is not listed either: the file browsers show only what can be opened. VitalStream's file
browser does not offer the folder above an allowed folder (its \ui{$\uparrow$} button is
disabled there).

\paragraph{Hidden entries.} Names that start with \code{.} are hidden:
\begin{itemize}
\item VitalStream's file browser lists them only when \ui{hidden} is ticked; a folder opened as
  one session leaves out the hidden entries and the symbolic links at every level; its index key
  leaves out the hidden entries, and counts a symbolic link to a file under the link's name, with
  its target's size and time;
\item \cmd{muse-py serve} and \cmd{dicom-py serve} never list or serve a hidden entry, nor
  anything inside a hidden folder;
\item the ECG and PACS pages pass over the hidden files and folders of a folder they are given.
\end{itemize}

\paragraph{The pages alone.} The ECG and PACS pages opened from disk or from the website read
nothing but what the user picks or drops, in the browser. They have no root: they are no server.
Opened with \code{\#src=URL} on a web server, they fetch that address with the page's own
credentials.

\subsection{The optional password}
\label{sec:viewers-alone-password}

All the viewers alone share one optional password, in one file:
\file{\textasciitilde/.config/vitalstream/auth.json} (VitalStream's setting \code{auth\_file},
\opt{-{}-auth-file FILE} of \cmd{serve}). It holds a salted scrypt hash of the password and a
random signing key, and is readable by its owner only. One password thus protects VitalStream,
\cmd{muse-py serve} and \cmd{dicom-py serve}.

\begin{lstlisting}[style=shell]
vitalstream-cpp --set-password   # or vitalstream-rs, -py, muse-py serve, dicom-py serve
vitalstream-cpp --clear-password   # remove it (for every viewer: the file is shared)
\end{lstlisting}

\begin{itemize}
\item \opt{-{}-set-password} asks for the password twice (at least 8 characters), writes the
  file, and exits: \code{Password set (FILE). Restart the viewer for it to take effect.} A viewer
  reads the file when it starts. A new password also makes a new signing key, which ends the
  earlier sign-ins.
\item With a password, the page asks for it (\ui{Password}, \ui{Sign in}); a wrong one is
  answered after a delay, with \code{Wrong password.} Each viewer, on each port, has a sign-in
  of its own (the same password).
\item In VitalStream, a sign-in lasts 30 days, and signing out ends it on the server too (a copy
  of the cookie no longer works). After wrong passwords from one address, the next one is
  checked only after 1, 3, 7, 15 \dots{} up to 30\,s; the C++ viewer answers an earlier try
  \code{429} (\cref{sec:viewer-api-signin}).
\item For \cmd{muse-py serve} and \cmd{dicom-py serve}, a sign-in lasts 30 days, until the server
  stops, or until a browser signs out, which ends every sign-in; after a wrong password the next
  try, from any browser, waits 1\,s, then 2, 4 \dots{} up to 60\,s.
\item \opt{-{}-no-auth} runs a viewer without asking for the password, even when one is set.
\item When it starts, a viewer says whether a password protects it; without one, it says that
  anyone who can connect to its port can read the files under its root.
\end{itemize}

\paragraph{Builds without password support.} A viewer built without password support
(\opt{-{}-no-password} of \cmd{install.sh}, \cref{tab:viewers-alone-build-options}) never asks for
a password, even when the shared file holds one, and \opt{-{}-set-password} and
\opt{-{}-clear-password} fail (\code{this viewer was built without password support}, or, for
\cmd{serve}, \code{this viewer was installed without password support}). Such a VitalStream viewer refuses to start where a
password is required (the setting \code{require\_password}, which Medorama sets for the viewer it
runs): it cannot be used as Medorama's viewer.

\subsection{What Medorama adds}
\label{sec:viewers-alone-medorama}

\begin{xltabular}{\linewidth}{@{}P{2.4cm}LL@{}}
\caption{A viewer alone, and inside Medorama.}
\label{tab:viewers-alone-medorama}\\
\toprule
{} & {\tabhead Alone} & {\tabhead Inside Medorama} \\
\midrule
\endfirsthead
\toprule
{} & {\tabhead Alone} & {\tabhead Inside Medorama} \\
\midrule
\endhead
\bottomrule
\endfoot
Install & VitalStream: \cmd{viewers/vitalstream/install.sh} (one viewer, a root it never
  leaves, by default the home folder, and an optional password). ECG and PACS: nothing (open the
  page), or \cmd{pip install} their Python package for \cmd{serve} &
  \cmd{tools/install-medorama.sh} builds Medorama and one VitalStream viewer (the whole file
  system as its root, with password support); the ECG and PACS pages come with Medorama \\
What it opens & The files and folders its user picks (VitalStream: under its root and its
  \code{data\_roots}; \cmd{muse-py serve}, \cmd{dicom-py serve}: under their \opt{-{}-root}) & What
  Medorama lets each user read, for a patient, from its catalog and connected databases \\
Who may use it & Whoever can reach its port; with a password set, only those who know it &
  Medorama's users, after Medorama's sign-in, with roles and permissions per user \\
Patient identity & Never shown by a viewer, whatever the file holds & The same; Medorama also
  de-identifies the files on its server (Privacy Mode, the default for every user but
  administrators) \\
\end{xltabular}

Patients, databases, users, roles, sign-in, audit and de-identification belong to Medorama only:
a viewer alone has none of them, beyond its optional password. How Medorama runs the viewers is
described in \cref{chap:viewers-in-medorama}, its users and Privacy Mode in
\cref{chap:users-privacy}.

\subsection{Patient data}
\label{sec:viewers-alone-privacy}

\begin{phibox}
The viewers show no patient identity. They read identifying fields of a file (a DICOM header's
name, a WFDB comment) only to group what belongs together, never to display them. The ECG and
PACS pages send nothing anywhere. VitalStream's server reads the recordings where they lie and
sends the page their samples, never a name, an identifier or a date of birth. A file's path and
its values may still identify a patient (bedside-monitor systems name files by unit, bed and
time): keep recordings, and the indexes made from them, on protected storage. In a clone, the
folder \file{sample\_data/} is the place for them; git ignores it.
\end{phibox}

% ==========================================================================================
\section{Opening records}
\label{sec:viewers-alone-opening}

The three viewers have the same two buttons, \ui{Open record\dots} and
\ui{Open multiple records\dots}:

\begin{xltabular}{\linewidth}{@{}P{2.3cm}LLL@{}}
\caption{The two buttons of the three viewers, and the limits on a folder.}
\label{tab:viewers-alone-opening}\\
\toprule
{} & {\tabhead VitalStream} & {\tabhead ECG} & {\tabhead PACS} \\
\midrule
\endfirsthead
\toprule
{} & {\tabhead VitalStream} & {\tabhead ECG} & {\tabhead PACS} \\
\midrule
\endhead
\bottomrule
\endfoot
\ui{Open record\dots} & One WFDB record (its \file{.hea}; a multi-segment record whole, a
  segment alone), an EDF or BDF file, or a file an adapter reads & One ECG: a WFDB record (its
  \file{.hea}, then its signal files), MUSE XML, DICOM, EDF, JSON, PNG, JPEG or PDF & One record:
  a study's or series' DICOM files, a video, a picture or a PDF \\
\ui{Open multiple records\dots} & A folder: every WFDB record in it and in its subfolders on one
  time axis (one session), or a WFDB collection; a summary first, then \ui{Open as one session}.
  EDF files are opened one at a time & A folder: every ECG in it and in its subfolders, EDF files
  included, stepped through one at a time (\ui{\textless}~\ui{\textgreater},
  \key{$\leftarrow$}~\key{$\rightarrow$}) & A folder: every study in it and in its subfolders,
  listed by series \\
Limits & 5000 header files, 20 folder levels, 10\,s to count them (settings) & 5000 ECGs & 5000
  files and 4\,GB \\
\end{xltabular}

VitalStream opens the files of the computer it runs on, through its own file browser. The ECG
and PACS pages open the files of the computer the browser runs on. When they are served by
\cmd{muse-py serve} or \cmd{dicom-py serve}, they also open those of the serving computer, whose
folders the button \ui{Server folders\dots} lists. The limits keep a folder opened by mistake
from overwhelming a viewer. When a limit is reached, the viewer says which one, and what to open
instead.

\paragraph{A WFDB header that leaves values out.} The ECG viewer applies WFDB's own defaults: a
sampling frequency of 250\,Hz, a gain of 200\,adu/mV (also for a gain of 0), the unit mV, and the
baseline at the ADC zero. VitalStream reads only what a header states: a record without a
sampling frequency is left out, and a signal without a gain is shown in \code{adu}.

% ==========================================================================================
\section{Light and dark}
\label{sec:viewers-alone-themes}

Every viewer has a light and a dark theme, switched with a button at the top right. In
VitalStream it is the half-filled circle; in the ECG and PACS viewers, it is a crescent moon in
the light theme and a sun in the dark one. The choice is remembered in the browser. Until a
choice is made, VitalStream and the ECG viewer follow the system's setting, and the PACS viewer
is dark. A host sets any of the three viewers' theme with \code{theme=light} or
\code{theme=dark} in the address (\cref{sec:viewers-alone-addresses}); the theme is then applied,
but not remembered as the viewer's own choice. Medorama passes its page's theme to every viewer
tab, when the tab opens and whenever the theme changes. The screenshots of this manual show the
dark theme, except the ECG viewer, which is shown in both (\cref{fig:ecg-viewer-main,fig:ecg-viewer-dark}).

% ==========================================================================================
\section{Addresses}
\label{sec:viewers-alone-addresses}

A viewer is told what to show by its address's fragment (the part after \code{\#}), so that
another program can open it at a file. Changing the fragment shows another file without
reloading the page.

\begin{xltabular}{\linewidth}{@{}P{1.9cm}P{4.0cm}L@{}}
\caption{The fragments of the viewers' addresses.}
\label{tab:viewers-alone-addresses}\\
\toprule
{\tabhead Viewer} & {\tabhead Fragment} & {\tabhead Means} \\
\midrule
\endfirsthead
\toprule
{\tabhead Viewer} & {\tabhead Fragment} & {\tabhead Means} \\
\midrule
\endhead
\bottomrule
\endfoot
ECG & \code{src=URL} & The file: a WFDB header (its signal files, and a multi-segment record's
  segments, are read beside it: \code{URL}'s folder followed by their names), a MUSE XML or
  DICOM ECG, an EDF or BDF file, a JSON ECG, a picture or a PDF \\
PACS & \code{src=URL} & One file \\
PACS & \code{list=URL} & A folder: \code{URL} answers \code{\{"files": ["IM0001.dcm", \dots]\}},
  and each file is read at \code{URL} followed by its name \\
PACS & \code{mm=\dots}, \code{fps=\dots} & The millimeters per pixel of pictures and videos; the
  frame rate of pictures (played as one loop) \\
ECG & \code{shift=\dots}, \code{cap90=1} & For a host that de-identifies in the page: the
  report's dates and the record's start moved back by that many seconds, ages above 89 years
  shown as 90+ (Medorama
  de-identifies on its server instead) \\
ECG, PACS & \code{name=\dots} & The name to show \\
All three & \code{theme=light} or \code{dark} & The colors, as the host's (VitalStream: applied
  without being saved as the viewer's own choice) \\
All three & \code{who=\dots}, \code{what=\dots}, \code{pos=\dots}, \code{nav=back,prev,next} &
  A host's context: whose file it is and what it is, shown in place of the file's name, with the
  buttons \ui{\textless{} Patient} and \ui{\textless}~\ui{\textgreater} that the viewer only reports
  to the page that hosts it (a message to the parent page, of the same origin); without them the
  file's name shows \\
ECG & \code{compact=1} & The paper alone, fitted to the frame (shown inside another page) \\
VitalStream & \code{path=\dots\&t0=\dots\&t1=\dots}, and the host's context above;
  \code{browse=FOLDER} & A recording on the server, at a time span (Unix seconds); the file
  browser, open at a folder \\
\end{xltabular}

The files are fetched with the page's own credentials (the same origin), so that a viewer served
by Medorama reads exactly what Medorama lets that user read. The requests VitalStream's page
makes to its server, and Medorama's addresses of the viewers, are described in
\cref{chap:viewer-api}.

% ==========================================================================================
\section{Tests}
\label{sec:viewers-alone-tests}

\begin{lstlisting}[style=shell]
make test-viewers      # the ECG and PACS pages' readers and drawing (Node's built-in test runner)
pytest viewers/ecg/python/muse/tests viewers/pacs/python/dicom/tests    # muse-py, dicom-py, serve
make test-python       # the same, with VitalStream's Python tests
\end{lstlisting}

\cmd{make test-viewers}, at the root of the repository, needs Node.js and no packages. It tests
the readers of the ECG viewer (its registry; WFDB single- and multi-segment, EDF, DICOM ECG and
JSON) and its paper drawing,
the DICOM reader of the PACS viewer (whose \file{dicom.js} the ECG viewer keeps a copy of, for
DICOM ECGs) and the MUSE reader. VitalStream has its own suites (\cref{chap:vitalstream}).
